\documentclass[11pt,letterpaper]{article}
\usepackage[T1]{fontenc}
\usepackage[utf8]{inputenc}
\usepackage{lmodern,microtype}
\usepackage[margin=0.79in,headheight=15pt]{geometry}
\usepackage{parskip,enumitem,longtable,tabularx,booktabs,array,xcolor,fancyhdr,hyperref}
\definecolor{navy}{HTML}{17314C}
\definecolor{teal}{HTML}{087F78}
\hypersetup{colorlinks=true,linkcolor=navy,urlcolor=teal}
\setlength{\emergencystretch}{3em}
\setlist[itemize]{leftmargin=*,nosep,topsep=5pt}
\setlist[enumerate]{leftmargin=*,itemsep=3pt}
\renewcommand{\arraystretch}{1.22}
\pagestyle{fancy}\fancyhf{}
\fancyhead[L]{\small Enterprise Data Modeling}\fancyhead[R]{\small Documentation Suite}
\fancyfoot[C]{\thepage}
\newcommand{\srcbook}{\textit{The Data Model Resource Book}, revised edition, Volume 2}
\newcommand{\ResponseLine}{\par\vspace{1.05em}\noindent\rule{\linewidth}{0.4pt}\par}
\newcommand{\ResponseArea}{\ResponseLine\ResponseLine\ResponseLine}
\begin{document}
\begin{center}
{\LARGE\bfseries\color{navy} Travel Industry Data Model Reference}\par
\vspace{0.25em}{\large Conceptual to logical reference | original adaptation}\par
\vspace{0.3em}{\small Original engineering study material \textbullet{} 2026 edition}
\end{center}
\vspace{0.8em}
\section{Scope and boundary}
Travel reference pattern based on the subject areas of Chapter 8 (printed pp. 355--392). This document is an original adaptation and does not reproduce the book's data diagrams, schema, or examples. The scenario assumed is: A travel platform coordinates a multi-supplier itinerary, rescheduling, cancelation and partial refunds.


\section{Candidate aggregates and invariants}
\begin{tabularx}{\linewidth}{@{}p{0.23\linewidth} p{0.38\linewidth} p{0.35\linewidth}@{}}
\toprule
\textbf{Aggregate / relation} & \textbf{Semantic constraint} & \textbf{Example} \\ 
\midrule
Reservation & Separate a held booking from fulfilled ticket & Seat held before payment \\ 
Travel segment & Represent ordered legs and connections & Air and rail segments in one itinerary \\ 
Traveler role & Assign passengers to bookings without duplicating person & Group reservation with individual travelers \\ 
Fare / agreement & Capture price rules and cancellation terms & Refund eligibility changes with fare class \\ 
Travel ledger & Track changes, fees, refunds and credits & Reissued ticket after a schedule change \\ 
\bottomrule
\end{tabularx}

\section{Core modeling decisions}
\begin{itemize}
\item Determine which concepts are independently identified and which are attributes or value objects.
\item Keep a business event and its associated status history separate when audit and replay matter.
\item Represent time-bounded participation with effective dating and a documented interval convention.
\item Avoid treating a generic subtype list as sufficient enforcement of business eligibility.
\item Keep product definition, commercial offer, agreement and operational fulfillment distinct when their lifecycles differ.
\item Record canonical terminology and aliases for cross-domain integration rather than duplicating physical storage.
\end{itemize}

\section{Reference constraints}
Use a unique natural-key constraint only if the domain guarantees uniqueness across the same scope and time interval. For many-to-many associations, use explicit intersection entities with role, provenance and dates. Never assume that an optional relationship can be replaced by an empty string. For temporal exclusion rules, specify transaction boundaries and concurrency tests.


\section{Integration and privacy}
\begin{itemize}
\item Document authoritative producer, permitted consumer, field-level classification and retention trigger for every shared entity.
\item Define an integration contract that maps identifiers, versions, enumerated states and recorded time.
\item Confirm that access decisions are enforced in the application and database; a logical diagram alone is not an access-control mechanism.
\item Track changes and consumer compatibility as schema versions, not as unexplained field replacements.
\end{itemize}

\section{Validation cases}
\begin{itemize}
\item How is a revised itinerary versioned?
\item Which supplier owns each segment?
\item What event triggers a refund obligation?
\item Create a counterexample that appears valid under the ERD but violates a narrative rule; propose a constraint or workflow test.
\item List at least two assumptions that require domain-owner approval.
\end{itemize}

\section{Related artifacts}
Pair this reference with Chapter 8 study guide, the model-review laboratory, the Data Dictionary Template/Blank, and the normalization and schema-design workbooks.


\vfill
\hrule
{\footnotesize\textbf{Provenance.} Len Silverston, The Data Model Resource Book, Revised Edition, Volume 2: A Library of Universal Data Models by Industry Types (Wiley, 2001). Chapter 8, pp. 355--392; all illustrative constraints are supplementary examples.}
\end{document}
